% ============================================================================
%  Declaración del uso de inteligencia artificial
% ============================================================================
\chapter*{Declaración del uso de inteligencia artificial}
\addcontentsline{toc}{chapter}{Declaración del uso de inteligencia artificial}
\markboth{Declaración del uso de inteligencia artificial}{Declaración del uso de inteligencia artificial}

En el desarrollo de este trabajo se emplearon herramientas de inteligencia
artificial generativa (modelos de lenguaje de gran escala) como asistente de
programación y de redacción. El uso se concentró en:

\begin{itemize}
  \item \textbf{Código de análisis y simulación.} Asistencia en la escritura y
  depuración de los cuadernos de Python y de los módulos de análisis de datos.
  Todos los resultados numéricos fueron ejecutados por el autor y verificados
  contra referencias independientes --- las trazas del PLC y los informes de
  servicio de la planta.
  \item \textbf{Revisión crítica de los análisis.} Contraste de hipótesis
  alternativas, detección de sesgos en los estimadores y verificación de
  consistencia entre documentos intermedios. Varias de las correcciones
  incorporadas en el Capítulo~\ref{cap:caracterizacion} --- en particular la del
  sesgo del estimador y la del diagnóstico por correlación instantánea ---
  surgieron de ese contraste.
  \item \textbf{Redacción y edición.} Asistencia en la organización y la
  redacción del documento en \LaTeX, a partir de material, resultados y
  decisiones del autor.
\end{itemize}

Las siguientes tareas \textbf{no} fueron delegadas: la formulación del
problema y de las hipótesis, la obtención y curaduría de los datos, las
decisiones de diseño de los controladores, la interpretación de los resultados
y las conclusiones. Ninguna cifra, figura o referencia de este documento
procede de una generación no verificada: los datos provienen del historiador
de la planta y de las trazas del PLC, y las referencias bibliográficas fueron
consultadas en su fuente.

El autor asume la responsabilidad plena por el contenido, la exactitud y la
integridad de este trabajo.
